
The purpose of this work is to study the concept of uniformly distributed sequences modulo 1 and to prove several criteria that guarantee that a sequence has this property, including Weyl's theorem. A sequence $(x_n)_{n\geq1}$ modulo 1 consists of observing only the fractional part of its terms $\partfrac{x_n}$. We will analyze how these values are distributed in the interval $[0,1)$, and in particular, when they are uniformly distributed over this interval. Over the course of this paper, we will develop effective tools to prove this property.

The paper is divided into four chapters. In the first of them, we will introduce the notion of real numbers modulo 1 and the quotient space $\mathbb{T}=\mathbb{R}/\mathbb{Z}$, known as the one-dimensional torus, which can be identified with the interval $[0,1)$ with its endpoints connected. Next, we will define the concept of an uniformly distributed sequence modulo 1 and prove some fundamental properties, accompanied by illustrative examples. We conclude this chapter by proving Weyl's uniform distribution theorem, which establishes the uniform distribution of the sequence $(n\alpha)$ for every irrational $\alpha$. To do so, we make use of the approximation of continuous functions by trigonometric polynomials. 

The second chapter focuses on developing the main tools to verify the uniform distribution of a sequence. First, we present the functional characterization, which states that a sequence is uniformly distributed modulo 1 if and only if, for any integrable Riemann function, the mean of the function evaluated over the values of the sequence converges to the integral of that function on $\mathbb{T}$. The second result is Weyl's criterion. This is one of the most important methods, if not the most important, to verify uniform distribution. It reduces the problem to the asymptotic calculation of sums of complex exponentials. As applications, the uniform distribution of $(n\alpha)$ is demonstrated again, and additional examples are studied, including sequences that are not uniformly distributed, $(\sin n)$.

In the third chapter, bounding techniques for sums of exponentials are explored. We will use Euler's summation formula. This formula allows us to express discrete sums in terms of integrals, which can be estimated with precision. Thus, by bounding them, we can verify uniform distribution using Weyl's criterion. We will show that the sequence $(a\log n)$ is not uniformly distributed modulo 1, while sequences of the form $(an^\sigma)$ with $0<\sigma<1$ are. The chapter concludes with the proof of Fejér's theorem, which guarantees the uniform distribution of a sequence $(g(n))$ assuming that the function $g$ and its derivative $g'$ satisfy certain growth, limit, and monotonicity conditions.

To conclude, we study Van der Corput's difference method. This method is based on Van der Corput's inequality, which relates sums of exponentials to their corresponding sums of differences. From this inequality, we demonstrate that a sequence is uniformly distributed if its sequences of differences are also uniformly distributed. As a primary application, we prove that every sequence defined by a polynomial with at least one non-constant irrational coefficient is uniformly distributed modulo 1. Furthermore, Van der Corput's difference method allows us to extend Fejér's theorem to more general functions, provided that their higher-order derivatives satisfy the appropriate growth conditions.